\scanpage{27}
dont la formation est compatible avec chgt % L1
de base (il suffit donc de le considérer % L2
sur $\mathrm{Spec}\,\mathbb Z$), avec \deleted{des} isom % L3
\[
T_\ell(\deleted{\mathbb Z}\mathbb Q(-1))\simeq\deleted{\mathbb Z}\mathbb Q_\ell(-1)={\underbrace{T_\ell(G_m)}_{\substack{\text{Module de Tate}\\\text{\inserted{inverse} de }G_m}}}^{-1} % L4
\]
et le cas échéant ($X$ sur $\mathbb Q$). % L5

On peut définir $\deleted{\mathbb Z}\mathbb Q(-1)$ comme % L6
$\mathrm R^2 f_*(\deleted{\mathbb Z}\mathbb 1_{P^1_X})$, où $f:P_X^1\longrightarrow X$ est la % L7
projection canonique. Posant % L8
\[
\deleted{\mathbb Z}\mathbb Q(-n)=\deleted{\mathbb Z}\mathbb Q(-1)^{\otimes n}\qquad\text{pour }n\geq0 % L9
\]
on peut prouver, à l'aide des % L10
axiomes déjà posés, que si $f:X-S$ % L11
est \inserted{lisse} projectif à fibres \inserted{géom.} connexes non vides, % L12
\deleted{partout} de dim relative $d$, alors % L13
\[
\mathbb R^{2d}f_*(\mathbb 1_X)\simeq\mathbb Z_S(-d) % L14
\]
et si on enlève l'hyp ``projectif'' mais % L15
seult $f$ quasi-projectif, on trouve encore % L16
\[
\mathbb R^{2d} f_!(\unsure{1|\mathbb 1}_X)\simeq\mathbb Z_S(-d) % L17
\]

% archive sheet number: 26
